\chapter{Introdução}\label{cap:introducao}

O futebol é um dos desportos com maior projeção mediática e económica, e a análise dos jogos apoia-se cada vez mais em dados. Clubes e operadores televisivos usam-nos para estudar o comportamento das equipas, preparar os jogos com os adversários e fundamentar decisões táticas. Grande parte desses dados é, no entanto, anotada à mão por analistas. A outra parte vem de sistemas de várias câmaras instalados no estádio, caros e ao alcance de poucas competições. Obter a mesma informação de forma automática, a partir do vídeo da transmissão televisiva, que existe para a maioria dos jogos profissionais, tornaria esta análise mais rápida e mais acessível.

É neste contexto que se enquadra a reconstrução do estado de jogo (\emph{Game State Reconstruction}, GSR), formalizada por \citet{somers2024soccernet}. O objetivo é passar do vídeo de uma única câmara para um mapa do campo. Nesse mapa, cada interveniente aparece com a sua posição, a equipa, a função (jogador, guarda-redes ou árbitro) e o número da camisola. Para isso, o sistema resolve várias tarefas em cadeia. Primeiro, deteta as pessoas em cada \emph{frame} e segue-as ao longo do vídeo, com a ajuda da reidentificação quando um jogador fica tapado ou sai de imagem. Depois, calibra a câmara para projetar cada jogador no campo. Por fim, identifica cada um.

É nesta última etapa que se centra este trabalho. Inclui duas tarefas, o reconhecimento do número da camisola e a atribuição da equipa, e, em ambas, a imagem dá pouco com que trabalhar. Numa transmissão televisiva, o número ocupa poucos píxeis, aparece muitas vezes tapado ou de lado e só se lê em parte dos \emph{frames}. As cores dos equipamentos mudam com a luz e nem sempre contrastam o suficiente. O próprio artigo que propõe a tarefa aponta o número como a maior dificuldade da sua \emph{baseline} \citep{somers2024soccernet}. São, no entanto, estas duas tarefas que dão uma identidade a cada interveniente. Sem elas, o mapa mostra posições, mas não diz quem ocupa cada uma.

Os trabalhos sobre estas duas tarefas deixam ainda uma pergunta em aberto. Os métodos académicos são avaliados isoladamente, muitas vezes em conjuntos de dados próprios e fora do sistema completo. As soluções apresentadas ao desafio SoccerNet são avaliadas no sistema completo, mas alteram vários módulos ao mesmo tempo, e o ganho que obtêm não se consegue repartir pelas alterações feitas. Fica por saber quanto vale, sozinha, uma melhor identificação dos jogadores.

É a essa pergunta que esta dissertação procura responder. Na \emph{baseline} do desafio, o número da camisola só fica certo em $15.3\%$ dos \emph{tracklets}, e a atribuição da equipa e do lado do campo falha em $28$ dos $159$ vídeos avaliados. Para melhorar estas duas tarefas, alteram-se quatro dos nove módulos do \emph{pipeline}. O reconhecimento do número da camisola é substituído, e os módulos da equipa e do lado do campo são reescritos. A reidentificação ganha um descritor da cor do equipamento, sem que o seu modelo mude. Os outros cinco módulos, a deteção, o \emph{tracking}, os pontos do campo, a calibração e a agregação por \emph{tracklet}, ficam como na \emph{baseline}. Assim, qualquer diferença de desempenho resulta das alterações introduzidas.

No número da camisola, o MMOCR, o sistema de deteção e reconhecimento de texto usado pela \emph{baseline}, dá lugar a um modelo multimodal de visão e linguagem, o Qwen2.5-VL-3B. O modelo lê apenas as costas do jogador, recortadas a partir da pose estimada. As leituras de cada \emph{tracklet} são depois combinadas numa decisão única, em que só contam aquelas em que o próprio modelo diz estar seguro. Na equipa, o agrupamento por aparência passa a verificar o próprio resultado e a recorrer à cor do equipamento quando a reidentificação não separa as duas equipas. O lado do campo passa a ser decidido por uma cascata de quatro critérios, o último dos quais usa o guarda-redes. As regras e os limiares dos dois módulos são escolhidos nos conjuntos de teste e de validação. Na equipa, cada ajuste só é aceite se corrigir vídeos que falhavam sem estragar nenhum que a versão anterior já acertava.

O sistema proposto e a \emph{baseline} são executados nos mesmos $159$ vídeos anotados, nas mesmas condições e sobre os mesmos \emph{tracklets}. Antes de os comparar, foi preciso corrigir um defeito de calibração na configuração distribuída da \emph{baseline}, que não faz parte das alterações deste trabalho. A exatidão do número sobe de $15.3\%$ para $42.6\%$ dos \emph{tracklets}, e a da equipa e do lado sobe de $83.5\%$ para $92.8\%$. No GS-HOTA, a métrica oficial do desafio, o ganho é de $16.57$ pontos no conjunto de teste. A métrica é ainda decomposta por atributo, o que permite ler em separado o contributo do número e o da equipa. O conjunto de treino não serviu para escolher nenhuma regra nem limiar, e nele o GS-HOTA sobe $13.85$ pontos, um ganho da mesma ordem do teste.

Tão importante como os ganhos é perceber até onde vão. No número da camisola, as experiências feitas mostram que, nos casos mais difíceis, os algarismos não se veem na imagem, e nenhuma transformação do recorte os recupera. Outra fração do erro vem do \emph{tracking}, que divide jogadores em vários \emph{tracklets} e troca identidades, fora dos módulos alterados. Na equipa, a inspeção visual dos $159$ vídeos mostra que a maioria das falhas que restam vem do classificador de função. Em vinte vídeos da validação, este classificador dá como árbitros os jogadores de amarelo. A mesma inspeção mostra ainda que, em cinco vídeos, a própria anotação tem os lados invertidos.

\section{Objetivos e Contributos}\label{sec:objetivos}

O objetivo geral desta dissertação é quantificar o contributo da identificação dos jogadores para a reconstrução do estado de jogo. Para isso, propõem-se novas soluções para as duas tarefas de identificação da \emph{baseline} do SoccerNet-GSR. São elas o reconhecimento do número da camisola e a atribuição da equipa e do lado do campo. Os restantes módulos do \emph{pipeline} mantêm-se inalterados, com exceção da reidentificação, à qual se acrescenta um descritor de cor.

Deste objetivo geral decorrem três objetivos específicos. O primeiro é melhorar o reconhecimento do número da camisola sem treinar modelos nem anotar dados, dentro das limitações do \emph{hardware} disponível. O segundo é tornar a atribuição da equipa e do lado do campo fiável nos vídeos em que a \emph{baseline} falha. O terceiro é avaliar estas alterações de modo que o ganho obtido lhes possa ser atribuído sem ambiguidade, tanto em cada módulo como no sistema completo.

Do trabalho realizado resultam quatro contributos. O primeiro é um módulo de reconhecimento do número da camisola que dispensa treino específico. A estimativa de pose isola a região dorsal do jogador, e um modelo multimodal de visão e linguagem de pequena dimensão lê o recorte. A decisão de cada \emph{tracklet} usa apenas as leituras com confiança suficiente (Secção~\ref{sec:met-jnr}). O segundo é um método de atribuição da equipa que verifica a qualidade do agrupamento e recorre à cor do equipamento quando as características de reidentificação não separam as equipas. O lado do campo é depois determinado por uma cascata de critérios (Secção~\ref{sec:met-equipa}). Já o terceiro é um protocolo experimental que isola o efeito das alterações propostas. A \emph{baseline} e o sistema proposto são avaliados nas mesmas condições, e o GS-HOTA é decomposto por atributo (Capítulo~\ref{cap:resultados}). Por fim, o quarto é uma quantificação da fiabilidade da regra da posição média, usada para atribuir o lado do campo, sobre a anotação completa dos 160 vídeos. Esta análise, que os trabalhos revistos não fazem, revelou também cinco vídeos com os lados invertidos na própria anotação (Secção~\ref{sec:res_lado_quantificacao}).


\section{Estrutura do Documento}\label{sec:estrutura}

A dissertação está organizada em cinco capítulos. O presente capítulo enquadra o problema, delimita o âmbito do trabalho e enuncia os seus objetivos e contributos. O Capítulo~\ref{cap:estado-arte} revê o estado da arte sobre o reconhecimento do número da camisola e a atribuição da equipa, incluindo as soluções apresentadas ao desafio SoccerNet. Termina com as métricas de avaliação e as lacunas que motivam este trabalho. O Capítulo~\ref{cap:metodologia} descreve as alterações introduzidas nos quatro módulos, o ambiente experimental e o procedimento de avaliação. O Capítulo~\ref{cap:resultados} apresenta os resultados de cada módulo e o seu efeito no GS-HOTA, sempre em comparação com a \emph{baseline}. Por fim, o Capítulo~\ref{cap:conclusao} sintetiza os principais contributos e aponta direções de trabalho futuro. No final do documento, o Apêndice A reúne os resultados por vídeo dos dois sistemas.
